% PAPER_STUDIO_PARAGRAPH:ABS-P1
Permutation sensitivity, the tendency of a language model's answer to shift when the positions of options in a multiple-choice prompt are reordered, threatens the reliability of model evaluation. Existing evidence documents this instability but offers no mechanistic account of when or why it occurs. This paper proposes a systematic study of permutation invariance that will characterize the conditions under which option order affects predictions. The planned evaluation will compare candidate explanations across a set of models and benchmark tasks, treating option-order variants as the controlled intervention. The central question is whether sensitivity reflects a shallow positional bias or a deeper sensitivity to contextual distribution. Without such an account, reported multiple-choice accuracies may conflate model competence with artifacts of prompt construction, and the proposed experiments are designed to resolve this ambiguity.
